\section{Capitolo 18: Alberi Binari e BST}

\begin{dizionario}{BST (Binary Search Tree)}
Albero dove per ogni nodo:
- Tutti i valori a \textbf{Sinistra} sono minori.
- Tutti i valori a \textbf{Destra} sono maggiori.
Permette ricerche veloci (come la ricerca binaria).
\end{dizionario}

\begin{esercizio}{Es. 36: Conteggio Nodi}
\textbf{TRACCIA:} Scrivere una funzione ricorsiva che conta quanti nodi ci sono in totale nell'albero.
\end{esercizio}

\begin{jsbox}
function contaNodi(root) {
    if (!root) return 0;
    return 1 + contaNodi(root.sx) + contaNodi(root.dx);
}
\end{jsbox}

\begin{esercizio}{Es. 37: Altezza Albero}
\textbf{TRACCIA:} Calcolare l'altezza dell'albero (il percorso più lungo dalla radice a una foglia).
\end{esercizio}

\begin{jsbox}
function altezza(root) {
    if (!root) return -1; // O 0, dipende dalla definizione. Qui -1 per albero vuoto.
    
    let hSx = altezza(root.sx);
    let hDx = altezza(root.dx);
    
    return 1 + Math.max(hSx, hDx);
}
\end{jsbox}

\begin{esercizio}{Es. 38: Verifica BST}
\textbf{TRACCIA:} Controllare se un albero è un BST valido (ogni nodo deve rispettare min/max ereditati dai genitori).
\end{esercizio}

\begin{jsbox}
function isBST(nodo, min = -Infinity, max = Infinity) {
    if (!nodo) return true;
    
    if (nodo.val <= min || nodo.val >= max) return false;
    
    return isBST(nodo.sx, min, nodo.val) && 
           isBST(nodo.dx, nodo.val, max);
}
\begin{esercizio}{Es. 44: Specchia Albero (Mirror)}
\textbf{TRACCIA:} Scrivere una funzione che inverte speculare (sx diventa dx) tutto l'albero.
\end{esercizio}

\begin{jsbox}
function invertTree(root) {
    if (!root) return null;
    
    // Scambio i figli
    let temp = root.sx;
    root.sx = root.dx;
    root.dx = temp;
    
    // Ricorsione
    invertTree(root.sx);
    invertTree(root.dx);
    
    return root;
}
\end{jsbox}

\begin{esercizio}{Es. 45: Albero Simmetrico}
\textbf{TRACCIA:} Verificare se un albero è speculare a se stesso (il sottoalbero sx è il mirror del dx).
\end{esercizio}

\begin{jsbox}
function isSymmetric(root) {
    if (!root) return true;
    
    function isMirror(n1, n2) {
        if (!n1 && !n2) return true; // Entrambi null
        if (!n1 || !n2) return false; // Uno solo null
        
        return (n1.val === n2.val) &&
               isMirror(n1.sx, n2.dx) &&
               isMirror(n1.dx, n2.sx);
    }
    
    return isMirror(root.sx, root.dx);
}
\end{jsbox}
